\documentclass[11pt,a4paper]{article}

% =========================================================
% Packages
% =========================================================
\usepackage[margin=2.5cm]{geometry}
\usepackage{graphicx}
\usepackage{booktabs}
\usepackage{tabularx}
\usepackage{array}
\usepackage{enumitem}
\usepackage{hyperref}
\usepackage{xcolor}
\usepackage{float}

% =========================================================
% Custom Commands
% =========================================================
\newcommand{\placeholder}[1]{\leavevmode{\color{gray}\itshape #1}}

% =========================================================
% Document
% =========================================================
\begin{document}

\section*{Document Information}

\begin{tabularx}{\textwidth}{>{\bfseries}l X}
Project / Feature & \placeholder{Nama fitur atau proyek} \\
Date & \placeholder{1 January 2026} \\
Role & \placeholder{Peran dalam proyek} \\
Status & \placeholder{Draft / Implemented / Verified / Completed} \\
Author & \placeholder{Irfan} \\
\end{tabularx}

\vspace{1em}

% =========================================================
\section{Problem}

\placeholder{
Jelaskan masalah yang ingin diselesaikan.
Apa kondisi awalnya?
Mengapa masalah ini perlu diselesaikan?
}

% =========================================================
\section{Requirements}

Tuliskan requirement utama sistem.

\begin{table}[H]
\centering
\begin{tabularx}{\textwidth}{l X l}
\toprule
\textbf{ID} & \textbf{Requirement} & \textbf{Priority} \\
\midrule

REQ-001 &
\placeholder{Sistem harus...} &
Must \\

REQ-002 &
\placeholder{Sistem sebaiknya...} &
Should \\

\bottomrule
\end{tabularx}
\end{table}

% =========================================================
\section{Constraints}

Batasan yang memengaruhi solusi:

\begin{itemize}
    \item \placeholder{Hardware constraint}
    \item \placeholder{Software constraint}
    \item \placeholder{Protocol / network constraint}
    \item \placeholder{Performance constraint}
    \item \placeholder{Time / cost constraint}
\end{itemize}

% =========================================================
\section{System Context}

Jelaskan posisi fitur ini di dalam sistem secara keseluruhan.

\placeholder{
Subsystem apa saja yang terlibat?
Siapa berkomunikasi dengan siapa?
}

% Contoh memasukkan gambar:
%
% \begin{figure}[H]
%     \centering
%     \includegraphics[width=0.8\textwidth]{figures/system-context.png}
%     \caption{System Context}
%     \label{fig:system-context}
% \end{figure}

% =========================================================
\section{Architecture}

Jelaskan architecture atau data flow solusi.

\placeholder{
Data berasal dari mana?
Diproses di mana?
Siapa menyimpan state?
Output dikirim ke mana?
}

% =========================================================
\section{Interfaces}

\begin{table}[H]
\centering
\begin{tabularx}{\textwidth}{l l l X}
\toprule
\textbf{Source} &
\textbf{Destination} &
\textbf{Protocol} &
\textbf{Description} \\
\midrule

\placeholder{Subsystem A} &
\placeholder{Subsystem B} &
\placeholder{CAN / UDP / API} &
\placeholder{Data atau command yang dikirim} \\

\bottomrule
\end{tabularx}
\end{table}

% =========================================================
\section{Alternatives and Trade-offs}

\subsection{Alternatives}

\begin{enumerate}
    \item \placeholder{Alternative A}
    \item \placeholder{Alternative B}
    \item \placeholder{Alternative C}
\end{enumerate}

\subsection{Trade-offs}

\begin{table}[H]
\centering
\begin{tabularx}{\textwidth}{l X X}
\toprule
\textbf{Criterion} &
\textbf{Alternative A} &
\textbf{Alternative B} \\
\midrule

Complexity &
\placeholder{Low / Medium / High} &
\placeholder{Low / Medium / High} \\

Reliability &
\placeholder{Notes} &
\placeholder{Notes} \\

Performance &
\placeholder{Notes} &
\placeholder{Notes} \\

Maintainability &
\placeholder{Notes} &
\placeholder{Notes} \\

\bottomrule
\end{tabularx}
\end{table}

% =========================================================
\section{Design Decision}

\textbf{Selected solution:}

\placeholder{Tuliskan solusi yang dipilih.}

\vspace{0.5em}

\textbf{Reasoning:}

\placeholder{
Kenapa solusi ini dipilih?
Trade-off apa yang diterima?
}

% =========================================================
\section{Implementation}

\placeholder{
Jelaskan implementasi secara singkat.
Tidak perlu menyalin source code.
Fokus pada bagaimana desain diwujudkan.
}

% =========================================================
\section{Verification}

Bagaimana memastikan implementasi memenuhi requirement?

\begin{table}[H]
\centering
\begin{tabularx}{\textwidth}{l l X l}
\toprule
\textbf{Test ID} &
\textbf{Requirement} &
\textbf{Test / Verification Method} &
\textbf{Result} \\
\midrule

TEST-001 &
REQ-001 &
\placeholder{Test procedure} &
\placeholder{PASS / FAIL} \\

\bottomrule
\end{tabularx}
\end{table}

% =========================================================
\section{Validation}

\placeholder{
Apakah solusi benar-benar menyelesaikan kebutuhan pengguna
atau kebutuhan operasional?

Bagaimana hal tersebut divalidasi?
}

% =========================================================
\section{Issues and Root Cause}

Jika ada masalah selama development:

\begin{description}[style=nextline]
    \item[Symptom]
    \placeholder{Apa yang terlihat gagal?}

    \item[Immediate Cause]
    \placeholder{Penyebab langsung}

    \item[Root Cause]
    \placeholder{Penyebab fundamental}

    \item[Corrective Action]
    \placeholder{Apa yang diperbaiki?}

    \item[Prevention]
    \placeholder{Bagaimana mencegah kejadian serupa?}
\end{description}

% =========================================================
\section{Result and Impact}

\placeholder{
Apa hasil akhirnya?

Contoh:
- meningkatkan reliability,
- mengurangi debugging time,
- mengurangi dependency,
- meningkatkan maintainability,
- meningkatkan observability.
}

% =========================================================
\section{Lessons Learned}

\begin{itemize}
    \item \placeholder{Pelajaran pertama}
    \item \placeholder{Pelajaran kedua}
    \item \placeholder{Pelajaran ketiga}
\end{itemize}

% =========================================================
\section{If I Redesign This}

\placeholder{
Jika sistem ini dibuat ulang dari awal,
apa yang akan dilakukan berbeda dan mengapa?
}

% =========================================================
\end{document}
